% FILENAME: paper/section-Gravity.tex

\section{Gravity: The Emergence of Spacetime}
\label{sec:Gravity}

\subsection{The Vacuum Condensate and Time Emergence}
\label{subsec:vacuum_condensate}

\todo{Author: Explain that space is not an empty stage; "empty" space mathematically has zero volume. State that a non-degenerate macroscopic metric ($\det(g) \neq 0$) requires a non-trivial gauge field condensate. Argue that a static universe also collapses to zero volume; thus, time evolution is geometrically required for space to exist.}

\begin{equation}
    \det(g_{\mu\nu}) \neq 0 \implies A^{(L)}_\mu \neq 0
    \label{eq:kinematicVacuumCondensate}
\end{equation}

\begin{listing}[H]
\begin{minted}{lean}
-- CGD.Quantum.Vacuum
theorem kinematicVacuumCondensate (u : Universe) (bulk : Set SpacetimePoint)
  [vol : CGD.Axioms.MacroscopicVolume u bulk] (x : SpacetimePoint) (hx : x ∈ bulk) :
  u.sd_sector.val ≠ (fun _ _ => 0)
\end{minted}
\caption{Proof that a non-zero spacetime volume forbids the trivial vacuum.}
\label{lean:kinematicVacuumCondensate}
\end{listing}

\begin{equation}
    \partial_0 A_i = 0 \implies \det(g_{\mu\nu}) = 0
    \label{eq:kinematicStaticUniverseDegeneracy}
\end{equation}

\begin{listing}[H]
\begin{minted}{lean}
-- CGD.Cosmology.BigBang
theorem kinematicStaticUniverseDegeneracy (u : Universe) :
  isStaticUniverse u →
  ∀ x, (urbantkeMetric (fun mu nu => curvatureSl2c u.sd_sector mu nu x)).det = 0
\end{minted}
\caption{A purely static universe collapses the emergent geometric volume to zero.}
\label{lean:kinematicStaticUniverseDegeneracy}
\end{listing}

\begin{equation}
    F_{\mu\nu} = \frac{1}{2} \epsilon_{\mu\nu}^{\phantom{\mu\nu}\rho\sigma} F_{\rho\sigma} \implies \text{Metric signature is degenerate or Euclidean}
    \label{eq:kinematicTimeEmergence}
\end{equation}

\begin{listing}[H]
\begin{minted}{lean}
-- CGD.Cosmology.TimeEmergence.Theorem
theorem kinematicTimeEmergence (u : Universe) (phaseRegion : Set SpacetimePoint)
  (h_symm : ∀ x ∈ phaseRegion, isFully4DSymmetric (fun mu nu => curvatureSl2c u.sd_sector mu nu x)) :
  ∀ x ∈ phaseRegion,
    ¬ isLorentzian (urbantkeMetric (fun m n => curvatureSl2c u.sd_sector m n x))
\end{minted}
\caption{A Lorentzian time axis is forbidden without spontaneous Euclidean symmetry breaking.}
\label{lean:kinematicTimeEmergence}
\end{listing}

\subsection{Unimodular CDJ Gravity}
\label{subsec:unimodular_cdj}

\todo{Author: Since we lack equations of motion, introduce the Capovilla-Dell-Jacobson (CDJ) algebraic constraint framework. Explain that the Cosmological Constant ($\Lambda$) is not manually inserted into an action, but emerges purely as an integration constant from the unimodular topological requirement.}

\begin{equation}
    \det(g_{\mu\nu}) = c \propto \Lambda^3 \neq 0
    \label{eq:macroscopicVacuumEmergence}
\end{equation}

\begin{listing}[H]
\begin{minted}{lean}
-- CGD.Gravity.DomainSeparation
theorem macroscopicVacuumEmergence 
  (u : Universe) (Λ : ℂ) (bulkVacuum : Set SpacetimePoint)
  (h_Λ_neq_zero : Λ ≠ 0) (k : ℂ) (hk : k ≠ 0)
  ... [Urbantke matrix adapter definition elided] ...
  (h_vacuum : isVacuumRegion bulkVacuum u Λ) :
  ∃ (c : ℂ), c ≠ 0 ∧ ∀ x y : bulkVacuum, 
    (cgdUnimodularMetricAdapter (fun m n => cgdAdjointCurvature u m n x.val)).det = c
\end{minted}
\caption{The Unimodular emergence of a constant macroscopic volume tied to $\Lambda$.}
\label{lean:macroscopicVacuumEmergence}
\end{listing}

\begin{equation}
    R_{\mu\nu} = 0
    \label{eq:macroscopicRicciFlatEmergence}
\end{equation}

\begin{listing}[H]
\begin{minted}{lean}
-- CGD.Gravity.DomainSeparation
theorem macroscopicRicciFlatEmergence
  (u : Universe) (urbantke_tetrad : TetradField)
  (metric_compat : ∀ x μ ν, metricFromTetrad urbantke_tetrad μ ν x = 
                           CGD.Gravity.urbantkeMetric (fun m n => curvatureSl2c u.sd_sector m n x) μ ν)
  ... [Capovilla-Dell-Jacobson constraints (Eq 2.2b/c) elided] ...
  (bulkVacuum : Set SpacetimePoint) (hOpen : IsOpen bulkVacuum) :
  ∀ x ∈ bulkVacuum, ∀ μ ν, ricciTensor (fun m n p => urbantkeMetric (fun a b => curvatureSl2c u.sd_sector a b p) m n) μ ν x = 0
\end{minted}
\caption{The derivation of Ricci-flat General Relativity via the Capovilla-Dell-Jacobson constraint.}
\label{lean:macroscopicRicciFlatEmergence}
\end{listing}

\subsection{Exact Lorentzian Macroscopic Witness}
\label{subsec:exact_lorentzian}

\todo{Author: Emphasize that this theorem is a \textbf{Capstone}. To prove this framework isn't just a vacuous Euclidean toy model, we have analytically verified an exact SU(2) gauge configuration that yields a strictly non-degenerate, negative determinant, Lorentzian spacetime metric at the origin.}

\begin{equation}
    \Sigma^{ab} - \frac{1}{3}\delta^{ab}\text{Tr}(\Sigma) = 0 \quad \text{and} \quad \det(g_{\mu\nu}) < 0
    \label{eq:dynamicExactLorentzianSolution}
\end{equation}

\begin{listing}[H]
\begin{minted}{lean}
-- CGD.Gravity.ExactSolutions.MainTheorem
theorem dynamicExactLorentzianSolution :
  ∃ (u : Universe) (x : SpacetimePoint), 
    (∑ μ ν ρ σ, epsilon4 μ ν ρ σ • (cgdAdjointCurvature u μ ν x * cgdAdjointCurvature u ρ σ x)) = 
    ((∑ μ ν ρ σ, epsilon4 μ ν ρ σ • (cgdAdjointCurvature u μ ν x * cgdAdjointCurvature u ρ σ x)).trace / 3) • 1 ∧
    isLorentzian (urbantkeMetric (fun m n => curvatureSl2c u.sd_sector m n x))
\end{minted}
\caption{Formal construction of an exact non-Abelian field yielding a macroscopic Lorentzian metric.}
\label{lean:dynamicExactLorentzianSolution}
\end{listing}

\subsection{Machian Motion of Topological Defects}
\label{subsec:machian_motion}

\todo{Author: Return to the physicist's critique: if the topological action's variation is zero, how do things move? Explain that the emergent Urbantke metric natively satisfies the Bianchi identities, satisfying the Geroch-Jang theorem premises. Conclude that defects are strictly forced to travel along timelike geodesics.}

\begin{equation}
    \nabla_\mu T^{\mu\nu} = 0 \implies x^\mu(\tau) \text{ is a timelike geodesic}
    \label{eq:machianTopologicalDefectMotion}
\end{equation}

\begin{listing}[H]
\begin{minted}{lean}
-- CGD.Gravity.GeodesicMotion
theorem machianTopologicalDefectMotion
  (isFutureDirectedTimelike : SpacetimePoint → (Fin 4 → ℝ) → Prop)
  (isTimelikeGeodesic : Set SpacetimePoint → Prop)
  (u : Universe)
  ... [Manifold smoothness, Contracted Bianchi, and Litlib.GerochJang axioms elided] ...
  (T_real : Fin 4 → Fin 4 → SpacetimePoint → ℝ) (γ : Set SpacetimePoint)
  (h_T_real_def : ∀ mu nu x, T_real mu nu x = (emergentStressEnergy (fun a b p => curvatureSl2c u.sd_sector a b p) mu nu x).re)
  (h_dominant_energy : ∀ x, (∃ mu nu, T_real mu nu x ≠ 0) → ∀ t t', 
    isFutureDirectedTimelike x t → isFutureDirectedTimelike x t' → ∑ a, ∑ b, T_real a b x * t a * t' b > 0)
  (h_localized : ∀ U : Set SpacetimePoint, IsOpen U → γ ⊆ U → closure {x | ∃ mu nu, T_real mu nu x ≠ 0} ⊆ U) :
  isTimelikeGeodesic γ
\end{minted}
\caption{The mathematical necessity of timelike geodesic motion derived from geometry.}
\label{lean:machianTopologicalDefectMotion}
\end{listing}
